\chapter{Checksums and streaming state}
\MARSChapterHeading{Checksums and streaming state}
\label{revision:crc}

The image file from the preceding chapter has a precise byte representation.
Suppose you copy it to another machine and want to notice an accidental change.
Comparing file lengths cannot detect a changed sample. A checksum compresses
information about the selected bytes into a value that can be recomputed and
compared with a stored value.

This chapter implements a cyclic redundancy check, or CRC, with a 32-bit state.
Its result is an error-detection value, not a unique identity: distinct messages
can have the same 32 bits. A CRC also does not establish who produced the bytes
or prevent someone from deliberately changing both data and checksum. Here its
purpose is to detect disagreement about the bytes a program received.

The complete \texttt{crc32.asm} in \texttt{revision/labs/files} supplies bit,
table and file routines. A second complete example, \texttt{crc32-stream.asm},
continues one message across two updates. You will derive the update rule,
explain why its state survives buffer boundaries, compare measured work and
finish with an independent checker for a bounded chunk. Bytewise reasoning,
logical shifts, function calls and the preceding chapter's file I/O are prerequisites.

\section{Specify the message and the CRC variant}

The name CRC-32 alone is insufficient to determine every convention. Our variant
uses reflected feedback mask hexadecimal \texttt{edb88320}, initial raw state
\texttt{ffffffff} and final XOR \texttt{ffffffff}. Reflected processing examines
low bits first and shifts the state to the right. Those choices must agree with
the writer of a stored checksum.

The check message is exactly nine ASCII bytes \texttt{123456789}.
It contains neither a newline nor a string terminator. Its reported checksum is
\texttt{cbf43926}; the state before the final XOR is \texttt{340bc6d9}.
Name these values differently in your notes. A raw state can continue a message;
a reported checksum has already undergone the final transformation.

The update interface receives source pointer in \texttt{\$a0}, selected count
from zero through 4096 in \texttt{\$a1}, and a zero-extended raw state in
\texttt{\$a2}. It returns the next zero-extended raw state in \texttt{\$v0}.
The table version also receives its prepared table pointer in \texttt{\$a3}.
An empty update returns the incoming state and reads neither source nor table.
The source may be unaligned; the word table must be four-byte aligned. Updates
write no caller data and preserve full saved registers, frame pointer and stack
pointer. Counts outside the stated range are outside this kernel contract.

Selecting the bytes is part of using the algorithm correctly. A checksum over
an allocation's padding answers a different question from a checksum over a
file's selected contents. The same applies to a text editor that silently adds
a line ending. Begin with the nine-byte fixture before investigating a mismatch
in a longer input.

\section{Follow one byte through eight feedback steps}

Let $r$ denote the current unsigned 32-bit state and $b$ the next unsigned byte.
First replace $r$ with $r\mathbin{\mathrm{xor}}b$. Then perform eight steps.
In each step, remember the old low bit, shift $r$ right logically by one, and
XOR the feedback mask into the shifted state if the remembered bit was one.
Remembering before shifting matters: the shift discards the bit that decides
whether feedback is required.

The feedback mask represents the selected generator polynomial in this bit
orientation. In polynomial arithmetic with coefficients zero and one, adding
and subtracting coefficients both act as XOR; there is no carry between bit
positions. The shift and conditional XOR implement a polynomial remainder
update. You can establish the program's behavior from the recurrence without
memorizing a long polynomial expression. Replacing XOR with ordinary integer
addition would introduce carries and implement a different operation.

For the first ASCII byte \texttt{31}, initial XOR gives \texttt{ffffffce}.
Its low bit is zero, so the first step produces \texttt{7fffffe7} without
feedback. That state's low bit is one. The next shift produces
\texttt{3ffffff3}; XOR with \texttt{edb88320} gives \texttt{d2477cd3}.
Continue the same rule for the remaining six steps.

\begin{figure}[htbp]
\centering
\begin{tabular}{r c l}
\toprule
Completed bit steps & Bit discarded in this step & Raw state, hexadecimal\\
\midrule
0 & --- & \texttt{ffffffce}\\
1 & 0 & \texttt{7fffffe7}\\
2 & 1 & \texttt{d2477cd3}\\
3 & 1 & \texttt{849b3d49}\\
4 & 1 & \texttt{aff51d84}\\
5 & 0 & \texttt{57fa8ec2}\\
6 & 0 & \texttt{2bfd4761}\\
7 & 1 & \texttt{f8462090}\\
8 & 0 & \texttt{7c231048}\\
\bottomrule
\end{tabular}
\caption{Eight reflected steps after XORing ASCII \texttt{1}, byte \texttt{31}, into initial state \texttt{ffffffff}. A discarded one selects feedback after the logical shift. The last row is the raw state after one byte, before any final XOR.}
\label{fig:revision-crc-first-byte}
\end{figure}


The worked inner loop expresses that decision directly:
\Needspace{9\baselineskip}
\MARSCodeLabel{Retain the old low bit before shifting}
\begin{minted}{mars-idea}
andi $t2, $v0, 1
dsrl $v0, $v0, 1
beqzc $t2, bit_no_feedback
nop
xor $v0, $v0, $t3
\end{minted}
Here \texttt{\$t3} contains the zero-extended feedback mask.
\texttt{bit\_no\_feedback} continues with the bit counter, which reaches zero
after eight steps. The outer loop advances the source and decrements its byte
count. The two loops count bit steps and source bytes respectively.

\section{Carry raw state across buffer boundaries}

Write $U(r,A)$ for the raw state after updating state $r$ with byte sequence $A$.
The next step depends only on the current state and next byte. Therefore,
for consecutive sequences $A$ and $B$,
\[
  U(r,A\mathbin{\|}B)=U(U(r,A),B),
\]
where $\|$ means concatenation. Splitting the source into buffers changes how
bytes become available; it does not change the state transition for a byte.
An empty buffer is an identity update: $U(r,\varnothing)=r$.

\begin{figure}[htbp]
  \centering
  \begin{tabular}{l l l}
    \toprule
    Stage               & New bytes      & Raw state, hexadecimal \\
    \midrule
    Initialize          & none           & \texttt{ffffffff}      \\
    First update        & \texttt{1234}  & \texttt{641c1f5c}      \\
    Continue that state & \texttt{56789} & \texttt{340bc6d9}      \\
    \bottomrule
  \end{tabular}
  \caption{A message split after four bytes has the same final raw state as one update over all nine bytes. Complementing \texttt{340bc6d9} produces the reported value \texttt{cbf43926}. The split does not introduce another initialization or finalization.}
  \label{fig:revision-crc-stream-state}
\end{figure}


The complete stream example updates the first four bytes with the bit routine,
retains the actual returned state, initializes a table, then updates the remaining
five bytes with the table routine. At \texttt{stream\_first}, \texttt{\$v0}
is \texttt{641c1f5c}. At \texttt{stream\_done}, it is \texttt{340bc6d9}.
Changing update methods between buffers is valid because both implement the
same raw transition. Resetting the state at the split would discard the prefix.

Finalization occurs once, after the complete message. Feeding the complemented
prefix value to the raw update would start the suffix with different bits.
For an empty complete message, no update changes initial \texttt{ffffffff};
its final XOR produces zero. That is compatible with an empty \emph{update}
returning its incoming raw state unchanged. The two statements concern different
interfaces.

\section{Replace eight steps with a prepared table}

There are only 256 possible low-byte values. Define $T[i]$ as the state obtained
by applying eight reflected steps to initial state $i$, without adding a byte.
The initializer builds all 256 entries from their indices, independent of old
memory contents. It writes exactly 1024 bytes and returns 256 entries.
For example, $T[0]=\texttt{00000000}$, $T[1]=\texttt{77073096}$,
$T[128]=\texttt{edb88320}$ and $T[255]=\texttt{2d02ef8d}$.

After XORing the next byte into $r$, separate that value into its low byte and
remaining high bits. During eight right shifts, none of those high bits reaches
the low-bit decision before the eighth shift has completed. Their contribution
is therefore simply $r\gg8$. The low byte contributes the eight-step result
stored in the table. Because shifts and this feedback recurrence distribute over
XOR, the contributions combine as
\[
  i=(r\mathbin{\mathrm{xor}}b)\mathbin{\mathrm{and}}255,
  \qquad r'=(r\gg8)\mathbin{\mathrm{xor}}T[i].
\]
The shift here is logical. The input byte affects the index and contributes no
bits to $r\gg8$ because it occupied only the low eight positions.

\Needspace{12\baselineskip}
\MARSCodeLabel{One table update for one source byte}
\begin{minted}{mars-idea}
lbu $t0, 0($a0)
xor $t0, $t0, $v0
andi $t0, $t0, 255
dsll $t0, $t0, 2
daddu $t0, $a3, $t0
lwu $t1, 0($t0)
dsrl $v0, $v0, 8
xor $v0, $v0, $t1
\end{minted}

The masked index selects an entry number. Multiplying it by four converts entries
to byte offsets before adding the full table address. An unscaled index could
select a misaligned word and would select the wrong bytes even on a machine
that permitted such a load. The table must be prepared before nonempty updates;
its mere allocation does not supply useful entries.

\section{Keep a 32-bit state correct in 64-bit registers}

The mathematical state has 32 bits, but MARS uses 64-bit general registers.
Our interface represents a state with zeroes in the upper 32 bits. A signed word
load of \texttt{edb88320} would instead extend its high bit through those positions.
A subsequent full-width logical shift would move the unwanted high bits into
the intended state. The initial word uses \texttt{lwu}, table entries use
\texttt{lwu}, and full-register constants explicitly establish zero extension.
A signed load followed by explicit extraction of its low 32 bits is an equivalent
implementation.

Unsigned byte loads matter too. Source byte \texttt{80} means 128; a signed byte
load would supply a negative full-register value unless it were corrected before
XOR. A 32-bit arithmetic right shift copies bit 31 and can introduce sign
extension, so it does not express the unsigned recurrence. A full 64-bit arithmetic
shift happens to agree with a logical shift while the upper half is zero, because
bit 63 is then zero; that alternative remains correct under this invariant.
The worked code uses logical shifts to state its intent directly. Trace a byte
above 127 and a state with bit 31 set when checking a narrowed operation.

Once a raw state is zero-extended, logical right shift keeps it so. XOR with the
zero-extended mask or a zero-extended table entry cannot create upper bits.
This gives a short invariant proving the representation after every step.
Preserving only the low word of a saved register is insufficient for the separate
calling convention, however. The tests seed saved registers with nonzero high
halves so that an accidentally narrowed restore is observable.

\section{Finish a file only after successful EOF}

\texttt{crc32\_file} receives a terminated path, scratch pointer, capacity from
one through 256 and method zero or one in \texttt{\$a0} through \texttt{\$a3}.
It returns the finalized checksum in \texttt{\$v0} and status in
\texttt{\$v1}: zero for success, one for open failure and two for read failure.
On failure, checksum is zero. Check status first; zero is also the legitimate
checksum of an empty file.

After a successful open, retain its actual descriptor and initialize raw state
once. For a positive read result $k$, update exactly those $k$ scratch bytes,
using the previous raw result as the next seed. Scratch capacity controls the
request, not the count of valid bytes after a short read. Continue until read
returns zero. A negative result is failure, including after a prefix has already
contributed to the state.

The file routine holds the descriptor, scratch pointer, capacity, method/helper
and raw state in saved registers across calls. It closes the acquired descriptor
on both EOF and read failure, and never closes a failed open result. The close
service uses \texttt{\$v0}, so the raw state must survive that service as well.
Only the success path complements it and publishes a checksum with status zero.
The scratch interval and the first 1024 table bytes are workspace; path and other
caller bytes remain unchanged.

A partial prefix has a perfectly defined raw state, but it does not stand for a
successfully read complete file. Returning that value with success would hide
an I/O failure. In the lab, an injected failure after five bytes verifies this
case separately from an ordinary missing path and an ordinary EOF. Actual host
runs independently compare both methods with Python's checksum, including the
PPM file produced in Chapter~\ref{revision:files}.

\section{Charge for preparation when comparing work}

Both update methods must agree before their costs are compared. The table loop
replaces repeated bit decisions with one indexed word load per byte. Its 1024-byte
lookup object consumes memory, must be initialized and creates a different access
pattern. Whether it is useful depends on how much input will reuse that preparation.

\begin{table}[htbp]
\centering
\begin{tabular}{l r r r r}
\toprule
Input & Bytes & Bit update & Table update & Build plus table\\
\midrule
Empty & 0 & 10 & 3 & 14351\\
One byte \texttt{80} & 1 & 73 & 16 & 14364\\
\texttt{123456789} & 9 & 519 & 120 & 14468\\
Zero bytes & 64 & 3632 & 835 & 15183\\
Values zero through 255 & 256 & 14668 & 3331 & 17679\\
Deterministic byte pattern & 4096 & 233206 & 53251 & 67599\\
\bottomrule
\end{tabular}
\caption{Completed simulated instructions for the supplied compact-branch routines, measured from direct function entry to return. Table construction costs 14348 instructions once. The final column charges that construction to one update; these are neither elapsed times nor hardware cycles. The last pattern uses sample $b_i=(37i+11)\bmod256$ at zero-based index $i$.}
\label{tab:revision-crc-work}
\end{table}


For nine bytes, the update alone falls from 519 instructions to 120, but building
its table first costs far more than either update. Even the 256-byte example uses
fewer instructions with the bit routine when construction is charged to that one
message. For 4096 bytes, construction plus table update uses 67599 instructions,
compared with 233206 for the bit method.

The table may instead be prepared once and reused across many raw updates. Add
construction once to their total, not once per buffer. The supplied file routine
builds its table once per table-method file call; it does not maintain a hidden
cache across file calls. A caller managing a persistent prepared table has a
different setup boundary and should say so explicitly.

Counts include explicit \texttt{nop} instructions and data-dependent feedback
branches. They do not establish elapsed-time or hardware-cycle ratios. Cache
effects, simulator work and file I/O need separate measurements. Reuse the public
benchmark's exact inputs when comparing a changed routine.

\section{Check a bounded chunk independently}

A later PNG decoder needs CRCs over selected fields rather than every byte in a
record. Complete \texttt{crc32\_check\_chunk} in \texttt{crc32-chunk.asm}.
Its selected buffer contains four big-endian length bytes, four type bytes,
$n$ payload bytes and four big-endian stored-checksum bytes. It accepts selected
counts from twelve through 4096 and requires exactly $n+12$ bytes.

\begin{figure}[htbp]
\centering
\begin{tabular}{l l l}
\toprule
Selected interval & Meaning & Included in this CRC?\\
\midrule
\texttt{[0,4)} & Big-endian payload length $n$ & No\\
\texttt{[4,8)} & Four type bytes & Yes\\
\texttt{[8,8+n)} & Payload bytes & Yes\\
\texttt{[8+n,12+n)} & Stored big-endian checksum & No\\
\bottomrule
\end{tabular}
\caption{The independent checker first proves $\mathit{selectedCount}=12+n$. It then computes the checksum over the contiguous type-plus-payload interval. For an empty \texttt{IEND} payload, the four type bytes alone produce \texttt{ae426082}.}
\label{fig:revision-crc-chunk-layout}
\end{figure}


Arguments are byte pointer, unsigned selected count and prepared word-aligned table
pointer in \texttt{\$a0/\$a1/\$a2}. Reject an out-of-range count before accessing
source or table. Otherwise read the four length bytes and prove the exact extent
before reading another field or accessing the table. Comparing $n$ with
$\mathit{selectedCount}-12$ avoids an overflowing addition when the encoded length is huge.

For a proved extent, compute the finalized CRC of type plus payload. Return it
in \texttt{\$v0}, with \texttt{\$v1}=0 if it matches the stored value or one
if it differs. Extent rejection returns zero and status two. Preserve caller
bytes and full saved state, make no console or file calls, and return through
the supplied continuation. The byte buffer may be unaligned; objects are disjoint.
You may reuse either completed update routine. The prepared table is available
if your design uses it.

An empty payload still has four covered type bytes. The twelve-byte \texttt{IEND}
example ends with stored bytes \texttt{ae 42 60 82}. Reading them as a native
little-endian word would compare against the wrong number. Covering the length,
excluding the type or covering the stored checksum would each compute a different
value. The lab makes these mistakes distinguishable with independent cases.

This checker does not validate type spelling, required chunk order, compressed
contents or a complete PNG. The unrestricted four type bytes are part of its
checksum input. Keep the function's evidence within that contract.

\section*{Sources and laboratory connection}
The raw transition, interfaces and ownership rules are implemented and checked
against the supplied MARS program. The
\href{https://www.w3.org/TR/2025/REC-png-3-20250624/}{PNG Third Edition Recommendation},
sections 5.3 and 5.5, specifies the chunk fields, CRC coverage, processing conventions
and stored byte order used by the independent exercise. Its size limits and error
precedence are our local interface. The lab follows actual states through a split
message, tests complete-file success versus failure and produces the bounded
checker without requiring a particular register allocation or frame.
